\documentclass[instructions.tex]{subfiles}
\begin{document}
 
\chapter{De la rechute dans le péché véniel.}
 
\primo Si vous tombez dans des péchés véniels par fragilité, ne vous en étonnez pas : personne n’est exempt de ces sortes de chutes : Le juste , dit l’Écriture , tombe sept fois , et se relève. Tombassiez-vous cent fois le jour, il faudrait vous relever cent fois. Ne vous découragez jamais pour ces sortes de fautes : le démon gagnerait plus par votre découragement que par vos chutes mêmes.
 
Mais, direz-vous , malgré mes résolutions , je fais toujours les mêmes fautes. Eh! vous seriez bien malheureux , si vous en faisiez de nouvelles ! Dieu en tirera sa gloire, si elles servent à vous rendre plus humble et plus vigilant. On doit présumer que ces chutes ne sont que de fragilité, 1. lorsqu’elles vous arrivent par surprise , et que, malgré le désir sincère de vous donner à Dieu , vous ne laissez pas que de tomber ; 2. lorsque vous en gémissez , que vous les craignez, et qu’elles vous causent d’abord des remords.
 
Ne vous éloignez pas des sacremens pour ces sortes de fautes ; ce serait une dangereuse illusion. Tâchez seulement de rendre vos fautes moins fréquentes et moins volontaires, détestez-les de plus en plus avec confusion : plus les tentations qui vous y portent vous affligent, moins il y a de consentement ; plus elles sont fréquentes, plus il y a d’occasions de mérites.
 
\secundo Si vous tombez sans remords dans des fautes vénielès, fréquemment ou de propos délibéré, ce n’est plus par fragilité, mais par malice. Oh ! que cet état est dangereux ! Ces infidélités habituelles et volontaires vous disposent à de grandes chutes. Celui , dit Jésus-Christ , qui est infidèle dans les petites choses , le sera dans les plus grandes\footnote{Qui in modico iniquus est, et in majori iniquus est. Luc. 16.}
 
Il ne faut cependant pas examiner scrupuleusement les motifs et les circonstances de ces sortes de chutes et de rechutes; examinez-les avec simplicité : ne cherchez point tant à connoître si elles sont commises par fragilité ou non , si elles sont grièves ou légères, si vous êtes dans un état de tiédeur ou de ferveur ; cette connoissance est réservée à Dieu. Humiliez-vous de ces fautes , gémissez-en , laissez-vous conduire par un prudent confesseur ; soumettez-lui vos lumières , et suivez ses avis. Sans cela , vous tomberez dans le scrupule et dans l’illusion.
 
\section{Résolutions}
 
\primo Je ne me découragerai plus de mes chutes dans le péché véniel, quelque nombreuses qu’elles puissent être. —

\secundo Mais je m’en humilierai devant Dieu , et je ferai tout ce qui dépendra de moi pour éviter ces sortes de péchés.—

\tertio Je ne conserverai plus d’affection pour aucun. 

\prayer{Mon Dieu, venez puissamment au secours de ma faiblesse , qui est extrême ; faites-moi la grâce de craindre tout ce qui vous déplaît , et de le fuir avec soin.}
 
\end{document}
